\section{Introduction}\label{sec:intro}

A \emph{magma} is a set equipped with a single binary operation, subject to
no axioms.  Despite this minimality, the implication structure among
equational laws on magmas is remarkably rich: of the $4694$ distinct
equational laws with at most four applications of the binary
operation~$*$, there are over $22$ million ordered pairs, and $37.12\%$
of them are true implications~\cite{bolan2025equational}.

The Equational Theories Project~\cite{bolan2025equational} settled the
complete implication graph through a combination of automated theorem
proving, finite counterexample search, and Lean~4 formalization.  Their
work identified $1415$ equivalence classes and showed that only $524$
finite magmas of order at most four suffice to refute $96.3\%$ of false
implications.  More recently, Berlioz and Melli\`es~\cite{berlioz2026latent}
introduced \emph{Stone pairings} $\langle E \mid A \rangle$, defined as
the probability that a random variable assignment satisfies equation~$E$
in a finite magma~$A$, and used principal component analysis of these
pairings to reveal geometric structure in the space of equational
theories.  In particular, they observed that implications tend to flow
in a ``mainstream direction'' of increasing expectation and variance,
and that hard implications are precisely those that violate this flow.

The SAIR Mathematics Distillation Challenge~\cite{sair2026challenge}
poses the equational implication problem in a knowledge-distillation
setting: participants must compress their understanding of the
implication graph into a text cheatsheet of at most $10$\,KB, which is
then provided to a language model that must predict the truth value of
each implication query.  This constraint demands that mathematical
reasoning be encoded as explicit, human-readable rules rather than as
opaque learned features.

\subsection{Prior work and limitations}
An initial rule-based predictor for this challenge achieved strong
performance on routine instances but only $52.75\%$ accuracy on the
hardest test split (denoted \textsc{Hard3}), barely exceeding the
$51.25\%$ always-false baseline.  The predictor employed fourteen
\emph{contradiction motifs}---syntactic patterns in the source and
target equations intended to signal non-implication---together with
structural rules based on variable counts, equation size, and
term shape.  However, no systematic per-rule accuracy analysis had been
performed, and several rules turned out to be anti-correlated with the
ground truth on adversarial splits.

\subsection{Contributions}
In this paper, we undertake a systematic diagnostic analysis of
structural prediction rules for equational implication.  Our
contributions are as follows.
\begin{itemize}[leftmargin=*,itemsep=2pt,topsep=2pt]
\item We prove \emph{projection collapse lemmas}
  (Theorem~\ref{thm:left-proj} and Theorem~\ref{thm:right-proj}),
  showing that any magma satisfying $x = x*y$ is a left-projection
  algebra and any magma satisfying $x = y*x$ is a right-projection
  algebra.  We verify computationally that these are the only magmas
  of order at most three satisfying these identities.

\item We introduce a diagnostic framework that treats each prediction
  rule as a binary classifier and measures its precision across
  multiple test splits.  Using this framework, we identify rules with
  accuracy below $35\%$ and either remove them or add necessary
  guard conditions (Propositions~\ref{prop:motif-guards}
  and~\ref{prop:structural-guards}).

\item We formulate new implication rules for \emph{bare} source
  equations with four or more distinct variables
  (Proposition~\ref{prop:bare-v4}), proving that such equations
  collapse the carrier set to at most one element under mild
  structural hypotheses.

\item We evaluate the refined rule set on four benchmark splits,
  demonstrating an overall accuracy improvement from $69.8\%$ to
  $79.1\%$ and a \textsc{Hard3} improvement from $52.75\%$ to $63.5\%$.
\end{itemize}

\subsection{Paper organization}
Section~\ref{sec:prelim} establishes notation and recalls the necessary
background on magmas and equational logic.  Section~\ref{sec:main}
contains our main results: the projection collapse lemmas, the
diagnostic framework, guard conditions for contradiction motifs,
and the new bare-source rules.  Section~\ref{sec:discussion} analyzes
the remaining error modes, discusses the hard2 regression, and connects
our findings to the Stone pairing framework of
Berlioz and Melli\`es~\cite{berlioz2026latent}.
Section~\ref{sec:conclusion} summarizes our contributions and
identifies directions for future work.
